% Generated draft: requires mathematical and source review.
\documentclass[11pt]{article}
\usepackage[margin=25mm]{geometry}
\usepackage{fontspec}
\setmainfont{Latin Modern Roman}
\usepackage{amsmath,amssymb,mathtools,unicode-math}
\setmathfont{Latin Modern Math}
\usepackage{enumitem,needspace,xurl,hyperref}
\setlength{\parindent}{0pt}
\setlength{\parskip}{5pt}
\setlength{\emergencystretch}{3em}
\newcommand{\N}{\mathbb{N}}
\newcommand{\Z}{\mathbb{Z}}
\newcommand{\Q}{\mathbb{Q}}
\newcommand{\R}{\mathbb{R}}
\newcommand{\C}{\mathbb{C}}
\newcommand{\F}{\mathbb{F}}
\newcommand{\A}{\mathbb{A}}
\newcommand{\PP}{\mathbb{P}}
\newcommand{\E}{\mathbb{E}}
\newcommand{\eps}{\varepsilon}
\newcommand{\dd}{\,\mathrm{d}}
\begin{document}
\section*{UnsolvedMath ID: 30005597}
\textbf{Source problem code:} catalogue entry
\subsection{Definitions and mathematical statement}
The supplied record groups three questions about magnetic Neumann eigenvalues. It does not supply all hypotheses needed to specify their operators and admissible domains uniquely. The following notation describes the apparent planar interpretation; the unresolved choices are stated explicitly.

For a planar domain \(\Omega\subset\R^2\) and a real magnetic potential \(A\), the formal magnetic Laplacian and its energy are
\[
L_A=(-i\nabla-A)^2,\qquad q_A[u]=\int_\Omega|(-i\nabla-A)u|^2\,\dd x.
\]
A magnetic Neumann realization is associated with a suitable form domain imposing no Dirichlet condition on the outer boundary. On a smooth boundary, its natural boundary condition is \(n\cdot(-i\nabla-A)u=0\), where \(n\) is the outward unit normal. Whenever the intended realization has discrete spectrum, write its eigenvalues, in nondecreasing order and counted with multiplicity, as \(\lambda_j^N\), \(j\geq1\). For a singular potential, the form domain and the realization at the singularity must also be specified.

\textbf{1. Second eigenvalue with Aharonov--Bohm flux.}
In a conventional planar normalization, a pole \(p=(p_1,p_2)\in\Omega\) and flux parameter \(\nu\in\R\) give
\[
A_{p,\nu}(x)=\nu\frac{(-(x_2-p_2),\,x_1-p_1)}{|x-p|^2},\qquad x\neq p,
\]
whose circulation around a positively oriented circle enclosing \(p\) is \(2\pi\nu\). Denote the intended magnetic Neumann eigenvalues on the punctured domain by \(\lambda_j^N(\Omega,p,\nu)\). The question is to determine the domain or domains maximizing \(\lambda_2^N\) under the constraints intended in the report. If the intended constraint is fixed area \(V>0\), its optimization form would be
\[
\sup_{(\Omega,p)\in\mathcal A_{V,\nu}}\lambda_2^N(\Omega,p,\nu),
\]
including whether this supremum is attained and, if so, by which configurations. Here \(\mathcal A_{V,\nu}\) is an \emph{unspecified} admissible class with \(|\Omega|=V\); the supplied passage does not determine this class, whether the pole is fixed or optimized, or the realization at the pole. Fixed area for this first question is suggested by the optimization context but is not explicitly stated in the quoted question. The quoted passage asks only what happens for the second eigenvalue and does not identify its proposed maximizing shape.

The preceding source fragment states that, for nonintegral \(\nu\), equality holds if and only if \(\Omega=\Omega^*\). The comparison inequality to which this equality refers, and the precise meaning and placement of \(\Omega^*\), are absent. This fragment therefore does not specify a second-eigenvalue conjecture.

\textbf{2. Constant-field ground-state comparison.}
In the planar constant-field interpretation, let
\[
A_\beta(x)=\frac{\beta}{2}(-x_2,x_1),\qquad \operatorname{curl}A_\beta=\beta,
\]
and let \(\lambda_1^N(\Omega,\beta)\) denote the lowest eigenvalue of the corresponding magnetic Neumann realization. Let \(D_V\) be a disk of area \(V\). The reconstructed statement asks whether, for each prescribed field strength \(\beta\) in the intended parameter range and each admissible bounded simply connected planar domain \(\Omega\) of area \(V>0\),
\[
\lambda_1^N(\Omega,\beta)\leq\lambda_1^N(D_V,\beta).
\]
The supplied record does not specify the domain regularity or the range of \(\beta\). Its reconstructed statement explicitly gives this direction of comparison, while the verbatim excerpt refers to an omitted inequality. The supplied literature assessment calls this a constant-curvature comparison, whereas the reconstructed statement calls it a constant-field comparison; the planar formulation above is consequently conditional on the latter interpretation. No equality characterization for this question is supplied.

\textbf{3. A pole approaching the boundary.}
The apparent intended setting is a fixed domain \(\Omega\), an Aharonov--Bohm flux \(\nu\), and interior poles \(p_\varepsilon\) satisfying
\[
\operatorname{dist}(p_\varepsilon,\partial\Omega)=\varepsilon,\qquad \varepsilon\downarrow0.
\]
The requested task is to determine the asymptotic behavior of
\[
\lambda_1^N(\Omega\setminus\{p_\varepsilon\},p_\varepsilon,\nu).
\]
Here the puncture is understood conditionally as a removed point, rather than a hole of positive radius. The source does not specify boundary regularity, the approach point or path, the flux restrictions, the singular realization, or the desired order of the asymptotic expansion. In particular, it does not establish that the answer depends only on \(\varepsilon\). These missing data prevent a unique fully quantified asymptotic statement.
\subsection{Short English statement}
The record asks for three results: characterize maximizers of a second magnetic Neumann eigenvalue with Aharonov--Bohm flux, establish a disk upper bound for the constant-field ground state among simply connected domains of fixed area, and determine the ground-state asymptotics as a flux pole approaches the boundary. The operator realization, admissible classes, and several parameter and boundary hypotheses remain unspecified. The displayed formulations retain these uncertainties and require human mathematical review.
\subsection{Sources}
\begin{itemize}
\item Supplied local UnsolvedMath registry record, permanent ID \(30005597\), title \emph{Three Magnetic-Laplacian Optimization Problems}, exact source problem code \texttt{catalogue entry}. The record attributes the questions to Oberwolfach Reports workshop discussions; no individual proposer is supplied. Its proposed year is \(2023\).
\item Supplied primary reference: \emph{Geometric Spectral Theory (2024)}, Oberwolfach Reports, pp.~2017--2106, DOI \texttt{10.4171/owr/2023/36}: \url{https://doi.org/10.4171/owr/2023/36}. The supplied literature-source label is ``Oberwolfach Report 36,'' with the same link. The citation's \(2024\) title date and its \(2023\) DOI year are retained as supplied.
\item Supplied literature assessment, checked on \texttt{2026-08-22}: the three questions are described as explicitly left open in the primary report, with no later complete resolution found. This is supplied dated curation, not an independently verified assessment. The background explicitly requests independent verification of time-sensitive status and cited claims.
\item The registry labels the problem \texttt{open} and its statement \texttt{corrected\_verified}, and reports that the reconstruction was checked against the official Oberwolfach report. These are supplied provenance claims; no independent source consultation or mathematical verification is asserted here.
\end{itemize}
\end{document}
